\documentclass{beamer}

\usepackage{cuhkbeamer}

\begin{document}

% ====================== Slide 1: TITLE SLIDE (REVISED) ======================
\begin{frame}
  \title{Airspace Access and Integration}
  \author{Chiu Yu KO}
  \institute{Low Altitude Economics}
  \date{Week 8}
  \titlepage
\end{frame}


% ====================== Slide 2: FORMAT & OBJECTIVES ======================
\begin{frame}[t]
\frametitle{Learning Objectives}
By the end of this class, you should be able to:
\begin{itemize}
\item \textbf{Explain} how safe operating capacity and demand create airspace congestion.
\item \textbf{Compare} queues, reservations, priority rules, peak pricing, and auctions as allocation mechanisms.
\item \textbf{Estimate} a simple delay externality and an operator's maximum affordable slot payment.
\item \textbf{Evaluate} how digital coordination and interoperability affect cost, safety, and scalability.
\item \textbf{Design} a phased airspace-integration roadmap with explicit evidence and decision gates.
\end{itemize}
\end{frame}

% ====================== Slide 3: ICEBREAKER ======================
\begin{frame}[t]
\frametitle{The 8:00 AM Capacity Conflict}
\textbf{Imagine that two operations request the same constrained corridor and time window.}
\begin{itemize}
\item One is a time-sensitive public-service delivery.
\item The other is a scheduled commercial passenger or cargo operation.
\end{itemize}
\vspace{0.3cm}
Only one request can be served safely within the available window.
\begin{itemize}
\item Should priority depend on arrival time, reservation, social value, price, or a predefined public-service rule?
\item \textbf{Quick poll (2 min):} Which allocation rule would you choose, and what problem might it create?
\end{itemize}
\end{frame}

% ====================== Session 1 Title Card ======================
\begin{frame}[c]
  \begin{center}
    \huge\color{cuhkpurple}\textbf{Session 1}\\
    \vspace{0.5cm}
    \Large Congestion Pricing, Slot Auctions \& Mechanism Design
  \end{center}
\end{frame}

% ====================== Slide 4: CONCEPT PRIMER 1 ======================
\begin{frame}[t]
\frametitle{Concept Primer: The Airspace Bottleneck}
Airspace capacity is not simply a fixed number of drones per minute.
\begin{itemize}
\item Safe capacity depends on route geometry, separation, aircraft performance, surveillance, communications, weather, contingency routes, human oversight, and interactions with other airspace users.
\item Vertiport capacity, charging, passenger processing, and available stands may become binding before the corridor itself.
\item Demand above reliable capacity creates delay, displacement, cancellation, or the need for another route or time.
\item Sandbox X includes controlled trials related to airspace management and simultaneous unmanned-aircraft operations, but does not establish a general published capacity for Victoria Harbour corridors.
\end{itemize}
\textbf{Managerial lesson:} Define the constrained resource and time interval before designing an allocation mechanism.
\end{frame}

% ====================== Slide 5: CONCEPT PRIMER 2 ======================
\begin{frame}[t]
\frametitle{Concept Primer: Congestion Externalities}
A congestion externality arises when one additional operation increases delay, risk-management effort, or disruption for other users without bearing the full resulting cost.
\begin{itemize}
\item A congestion charge aims to reflect the additional cost imposed on others.
\item The relevant cost may include aircraft operating cost, passenger or cargo delay, schedule disruption, and additional coordination effort.
\item Pricing is only one response. Reservations, quotas, priority rules, rerouting, and capacity investment may also be appropriate.
\item Some public-service missions should be governed by explicit priority rules rather than willingness to pay.
\end{itemize}
\end{frame}

% ====================== Slide 6: STANDALONE FORMULA ======================
\begin{frame}[t]
\frametitle{Applied Tool: A Simple Marginal Congestion Cost}
\begin{block}{Illustrative Calculation}
\[
\begin{aligned}
\text{Marginal Congestion Cost} = {} & \sum_j \left(\text{Additional Delay}_j \times \text{Value of Delay}_j\right) \\
& {} + \text{Other Incremental Costs}
\end{aligned}
\]
\end{block}
\begin{itemize}
\item Sum the incremental delay imposed on all affected users, not only the next aircraft.
\item Values may differ across passengers, cargo, aircraft, public services, and schedule commitments.
\item The theoretically efficient charge equals marginal external cost, but that cost is difficult to observe precisely in real time.
\item A practical peak charge may therefore use estimated bands and periodic review.
\end{itemize}
\end{frame}

% ====================== Slide 7: BACK-OF-ENVELOPE (CONGESTION) ======================
\begin{frame}[t,shrink=11]
\frametitle{Back-of-the-Envelope: Estimating a Delay Externality}
\textbf{Scenario:} One additional operation delays three other flights by five minutes each.
\begin{columns}[T]
\begin{column}{0.48\textwidth}\begin{block}{Affected Users}
\begin{itemize}
\item Flight 1 delay value: HK\$100/min
\item Flight 2 delay value: HK\$60/min
\item Flight 3 delay value: HK\$40/min
\item Total value: HK\$200/min
\end{itemize}\end{block}\end{column}
\begin{column}{0.48\textwidth}\begin{block}{Estimated External Cost}
\begin{itemize}
\item Delay: 5 minutes
\item External delay cost: $5\times200$
\item \textbf{HK\$1,000}
\item A peak charge near this amount may encourage a flexible operation to move, but its response is not guaranteed.
\end{itemize}\end{block}\end{column}
\end{columns}
\textit{Takeaway: A congestion charge should reflect the total incremental cost imposed on others, while priority public-service operations require separate rules.}
\vfill{\tiny Illustrative assumptions; excludes safety, environmental, and schedule-recovery costs.}
\end{frame}

% ====================== Slide 8: SESSION 2 TITLE ======================
\begin{frame}[c]
  \begin{center}
    \huge\color{cuhkpurple}\textbf{Session 2}\\
    \vspace{0.5cm}
    \Large Slot Auctions, Mechanisms \& Network Externalities
  \end{center}
\end{frame}

% ====================== Slide 9: CONCEPT PRIMER 3 ======================
\begin{frame}[t]
\frametitle{Concept Primer: Mechanism Design and Slot Allocation}
Mechanism design asks how rules shape participants' information, incentives, and outcomes.
\begin{itemize}
\item \textbf{Administrative rules:} schedules, rotation, first-come-first-served, historical rights, or public-service priority.
\item \textbf{Posted prices:} predictable peak and off-peak charges.
\item \textbf{Auctions:} participants submit bids for scarce time windows or route access.
\item \textbf{Package allocation:} linked departure, corridor, arrival, and charging reservations may need to be allocated together.
\end{itemize}
\textbf{Design criteria:} safety, efficiency, fairness, competition, transparency, reliability, simplicity, and access for new entrants.
\textit{Failure to obtain a linked slot should lead to rejection or rescheduling, not an unsafe flight.}
\end{frame}

% ====================== Slide 10: STANDALONE FORMULA ======================
\begin{frame}[t]
\frametitle{Applied Tool: Maximum Affordable Slot Payment}
\begin{block}{A Simple Operator Screen}
\[
\text{Maximum Affordable Payment}
=
\text{Incremental Revenue}
-\text{Incremental Operating Cost}
-\text{Required Contribution}
\]
\end{block}
\begin{itemize}
\item Use incremental values associated with obtaining the slot, not total company revenue and OPEX.
\item The required contribution may cover fixed cost, risk, or a target return.
\item The rational bid also depends on auction rules, expected rival bids, alternative slots, and the risk that linked resources are unavailable.
\item A high bid does not automatically imply high social value or justify exclusion of smaller operators.
\end{itemize}
\end{frame}

% ====================== Slide 11: BACK-OF-ENVELOPE (AUCTIONS) ======================
\begin{frame}[t,shrink=11]
\frametitle{Back-of-the-Envelope: Comparing Affordable Bids}
\textbf{Scenario:} Two operators value the same commercial slot.
\begin{columns}[T]
\begin{column}{0.48\textwidth}\begin{block}{Operator A}
\begin{itemize}
\item Incremental revenue: HK\$2,000
\item Incremental cost: HK\$1,200
\item Required contribution: HK\$300
\item \textbf{Maximum affordable payment: HK\$500}
\end{itemize}\end{block}\end{column}
\begin{column}{0.48\textwidth}\begin{block}{Operator B}
\begin{itemize}
\item Incremental revenue: HK\$2,000
\item Incremental cost: HK\$600
\item Required contribution: HK\$300
\item \textbf{Maximum affordable payment: HK\$1,100}
\end{itemize}\end{block}\end{column}
\end{columns}
\textit{Takeaway: Operator B can afford a higher bid under these assumptions. The actual payment depends on the auction design and competing bids, not necessarily HK\$501.}
\vfill{\tiny Illustrative assumptions; commercial slots remain subject to safety, priority, and competition rules.}
\end{frame}

% ====================== Slide 12: CONCEPT PRIMER 4 ======================
\begin{frame}[t]
\frametitle{Concept Primer: Network Effects and Shared Information}
A network effect exists when participation by additional compatible users increases the value of the network to existing users.
\begin{itemize}
\item Common interfaces and timely information may improve route planning, conflict detection, situational awareness, and coordination.
\item More participants can also increase congestion, cybersecurity exposure, monitoring needs, and system complexity.
\item Shared data does not automatically make a system safer. Data quality, latency, coverage, governance, and operational use matter.
\item Some benefits are network effects; others are economies of scale or standardization benefits. They should not be treated as identical.
\end{itemize}
\end{frame}


% ====================== Slide 13: STANDALONE FORMULA ======================
\begin{frame}[t]
\frametitle{Applied Tool: Coordination-Cost Savings}
\begin{block}{A Simple Comparison}
\[
\begin{aligned}
\text{Daily Net Saving} = {} & \text{Cost of the Less-Integrated Process} \\
& {} - \text{Cost of the Interoperable Process}
\end{aligned}
\]
\end{block}
\begin{itemize}
\item Each process should include fixed systems, usage fees, communications, human oversight, cybersecurity, integration, certification, and exception handling.
\item Interoperability can reduce duplicate interfaces and repeated manual data entry.
\item It may also create migration cost, common dependencies, governance disputes, and switching risks.
\item Avoid assuming that API or digital-services cost is near zero.
\end{itemize}
\end{frame}

% ====================== Slide 14: BACK-OF-ENVELOPE (U-SPACE) ======================
\begin{frame}[t,shrink=11]
\frametitle{Back-of-the-Envelope: Two Coordination Processes}
\textbf{Assumptions:} 1,000 operations/day and 30 operating days/month.
\begin{columns}[T]
\begin{column}{0.48\textwidth}\begin{block}{Less-Integrated Process}
\begin{itemize}
\item Manual and system cost: HK\$30/operation
\item Daily cost: HK\$30,000
\item \textbf{Monthly cost: HK\$900,000}
\end{itemize}\end{block}\end{column}
\begin{column}{0.48\textwidth}\begin{block}{Interoperable Digital Process}
\begin{itemize}
\item Fixed system and oversight: HK\$12,000/day
\item Variable cost: HK\$3/operation
\item Daily cost: HK\$15,000
\item \textbf{Monthly cost: HK\$450,000}
\end{itemize}\end{block}\end{column}
\end{columns}
\textit{Takeaway: The interoperable process saves HK\$450,000 per month under these assumptions, but only if it achieves the required reliability, security, and service quality.}
\vfill{\tiny Illustrative assumptions; not a Hong Kong UTM or European U-space cost estimate.}
\end{frame}

% ====================== Slide 15: SESSION 3 TITLE ======================
\begin{frame}[c]
  \begin{center}
    \huge\color{cuhkpurple}\textbf{Session 3}\\
    \vspace{0.5cm}
    \Large Interoperability, Roadmaps \& GBA Integration
  \end{center}
\end{frame}

% ====================== Slide 16: CONCEPT PRIMER - INTEROPERABILITY ======================
\begin{frame}[t]
\frametitle{Concept Primer: Interoperability and Governance}
Interoperability is the ability of authorized systems and organizations to exchange, interpret, and use information consistently.
\begin{itemize}
\item It may rely on common data models, interfaces, identity and access controls, service levels, and conformance testing.
\item Shared standards can have public-good characteristics, but the systems, operations, cybersecurity, and assurance remain costly.
\item Open standards do not require unrestricted public access to operational data.
\item Governance must address data minimization, confidentiality, liability, competition, resilience, vendor dependence, and incident response.
\end{itemize}
\textbf{Managerial lesson:} Technical compatibility is necessary but not sufficient; institutions and incentives must also align.
\end{frame}

% ====================== Slide 17: APPLIED TOOL - INTEROPERABILITY SAVINGS ======================


% ====================== Slide 18: BACK-OF-THE-ENVELOPE - INTEROPERABILITY ======================


% ====================== Slide 19: CONCEPT PRIMER - AIRSPACE INTEGRATION ROADMAP ======================
\begin{frame}[t]
\frametitle{Concept Primer: Airspace-Integration Roadmap}
Integration should proceed through evidence-based stages rather than fixed promises of full scale by a particular year.
\begin{itemize}
\item \textbf{Stage 1: Defined trials} with limited routes, aircraft, times, users, and contingency procedures.
\item \textbf{Stage 2: Extended operations} with more conditions, repeated performance, interoperable services, and stakeholder feedback.
\item \textbf{Stage 3: Scaled network} only after safety, reliability, capacity, governance, and commercial thresholds are met.
\item Cross-boundary development also requires coordination on customs, immigration, quarantine, spectrum, data, security, liability, and emergency response.
\end{itemize}
\end{frame}

% ====================== Slide 20: APPLIED TOOL - ROADMAP VALUATION (REVISED) ======================
\begin{frame}[t]
\frametitle{Applied Tool: Phased Roadmap Value}
\begin{block}{A Simple NPV Framework}
\[
\text{Roadmap NPV} = \sum_{t=0}^{T} \frac{
  \begin{aligned}
    & \text{Incremental Cash Inflow}_t \\
    {} - {} & \text{Incremental Cash Outflow}_t
  \end{aligned}
}{(1+r)^t}
\]
\end{block}
\begin{itemize}
\item Include implementation cost in the year in which it occurs, rather than subtracting all costs outside the discounting formula.
\item Cash flows should reflect probability of reaching each phase, operating limits, demand, cost, and possible abandonment.
\item Decision gates preserve flexibility to pause, redesign, or expand when evidence changes.
\item A positive NPV under assumptions does not prove that the wider market is viable.
\end{itemize}
\end{frame}
% ====================== Slide 21: BACK-OF-THE-ENVELOPE - ROADMAP (REVISED) ======================
\begin{frame}[t,shrink=12]
\frametitle{Back-of-the-Envelope: A Three-Year Phased Roadmap}
\textbf{Assumptions:} Discount rate = 10\%. Net cash flows already combine phase revenue and implementation or operating cost.
\begin{columns}[T]
\begin{column}{0.48\textwidth}\begin{block}{Cash Flows}
\begin{itemize}
\item Year 0 trial investment: $-$HK\$50M
\item Year 1 net cash flow: HK\$10M
\item Year 2 net cash flow: HK\$40M
\item Year 3 net cash flow: HK\$50M
\end{itemize}\end{block}\end{column}
\begin{column}{0.48\textwidth}\begin{block}{Present Values}
\begin{itemize}
\item Year 1: HK\$9.1M
\item Year 2: HK\$33.1M
\item Year 3: HK\$37.6M
\item Total PV of future net cash flows: HK\$79.8M
\item \textbf{Roadmap NPV: HK\$29.8M}
\end{itemize}\end{block}\end{column}
\end{columns}
\textit{Takeaway: The roadmap is attractive under these assumed cash flows. Probabilities, decision gates, terminal value, and cross-boundary dependencies still require analysis.}
\vfill{\tiny Illustrative assumptions; not a forecast of Hong Kong-GBA operations.}
\end{frame}

% ====================== Slide 19: EXERCISE INSTRUCTIONS ======================
\begin{frame}[t]
\frametitle{In-Class Exercise: Allocating a Scarce Slot}
\textbf{Scenario:} Disruption reduces three planned corridors to one available commercial slot.
\begin{itemize}
\item \textbf{Roles:} Commercial operators, public-service customer, regulator or airspace manager, and community representative
\item \textbf{Task:}
\begin{enumerate}
\item calculate each commercial operator's maximum affordable payment;
\item choose an allocation mechanism and explain its objective;
\item apply a predefined public-service priority rule if triggered;
\item state how cancelled or displaced operations are treated; and
\item assess effects on efficiency, fairness, competition, and daily breakeven.
\end{enumerate}
\item \textbf{Timing:} 15 minutes plus 5-minute discussion.
\end{itemize}
\end{frame}

% ====================== Slide 20: EXERCISE TEMPLATE ======================
\begin{frame}[t]
\frametitle{Exercise Template: Allocation Ledger}

\begin{center}\small
\setlength{\tabcolsep}{5pt}
\begin{tabular}{l|c|c|c}
\textbf{Request} & \textbf{\shortstack{Incremental\\revenue}} & \textbf{\shortstack{Incremental\\cost}} & \textbf{\shortstack{Required\\contribution}}\\
\hline
Cargo operator & HK\$800 & HK\$400 & HK\$100\\
Passenger operator & HK\$2,000 & HK\$600 & HK\$400\\
Priority public service & \multicolumn{3}{c}{Apply predefined eligibility and dispatch rule}\\
\end{tabular}
\end{center}

\vspace{0.2cm}
\begin{itemize}
\item Cargo maximum affordable payment: HK\$300.
\item Passenger maximum affordable payment: HK\$1,000.
\item If priority is triggered, record who bears the displacement cost and whether compensation, rescheduling, or reserved capacity applies.
\end{itemize}
\end{frame}

% ====================== Slide 21: KEY TAKEAWAYS ======================
\begin{frame}[t]
\frametitle{Key Takeaways}
\begin{itemize}
\item Airspace capacity depends on the complete operating system and cannot be reduced to an unsupported fixed number of drones per minute.
\item Congestion pricing estimates external delay cost, while queues, reservations, priority rules, and auctions serve different objectives.
\item Auctions can reveal willingness to pay but require safeguards for safety, competition, linked resources, new entrants, and public-service access.
\item Interoperability may reduce duplicate coordination and support scale, but does not make digital services free or automatically safe.
\item Network participation can create benefits and costs; data quality, cybersecurity, governance, and congestion remain essential.
\item A credible roadmap links each expansion in capacity or operating privilege to evidence, review, and a financially consistent decision gate.
\end{itemize}
\end{frame}

% ====================== Slide 22: ASSIGNMENT ======================
\begin{frame}[t]
\frametitle{Week 8 Assignment Briefing}
\textbf{Airspace Allocation and Congestion Memo (maximum 1,200 words)}
\begin{itemize}
\item \textbf{Due:} Beginning of Week 9 (group assignment)
\item \textbf{Task:} Design an allocation mechanism for a constrained Hong Kong corridor or facility.
\item \textbf{Requirements:}
\begin{enumerate}
\item define the constrained resource, safe capacity, time interval, and demand scenario;
\item compare at least two mechanisms, such as queueing, reservations, peak pricing, or auction;
\item estimate a delay externality or maximum affordable payment using transparent assumptions;
\item specify eligibility and treatment for priority public-service operations;
\item model one weather or capacity shock and identify who bears displacement cost;
\item explain the required digital information, interoperability, cybersecurity, and governance arrangements; and
\item propose a phased implementation and evaluation plan.
\end{enumerate}
\item \textbf{Grading focus:} Economic reasoning (35\%), operational realism (25\%), welfare and competition analysis (20\%), governance and roadmap (10\%), professionalism (10\%).
\end{itemize}
\end{frame}

% ====================== Slide 23: PREVIEW WEEK 9 ======================
\begin{frame}[t]
  \frametitle{Preview of Week 9}
  
  \begin{center}
      \Large \textbf{Economic and Social Impact of LAE}
  \end{center}
  
  \vspace{1cm}
  \textbf{Preview question for next week:} \\
  \textit{We've spent 8 weeks figuring out how companies make money... but does the Low Altitude Economy actually make life better for the everyday citizens of Hong Kong?}
\end{frame}

% ====================== Slide 24: THANK YOU ======================
\begin{frame}[t]
  \frametitle{Thank You + Q\&A}
  
  \begin{center}
      \textbf{Next Week Preview:} \\
      \Large The Big Picture: Jobs, GDP, and the Environment
  \end{center}
  
  \vspace{0.5cm}
  \textbf{Pre-Class Readings Reminder (for Week 9):}
  \begin{itemize}
      \item \textit{NASA UAM Market Study (Network design \& airspace)}
      \item \textit{EASA Easy Access Rules for U-space (Interoperability)}
  \end{itemize}
  
  \vspace{0.5cm}
  \begin{center}
      \textbf{Questions?} \\
      \small (Final 15 minutes reserved for extended Q\&A)
  \end{center}
\end{frame}

\end{document}